\begin{summarybox}
\textbf{一句话理解：}椭圆上点 $P$ 看两个焦点的张角，
在\textbf{短轴端点}最大；因为此时焦点三角形是\textbf{等腰三角形}。

\medskip
\textbf{三个必须记牢的关系：}
\[
PF_1+PF_2=2a,\qquad F_1F_2=2c,\qquad e=\frac ca.
\]
\[
\boxed{\sin\frac\theta2\le e},
\qquad
\boxed{\theta_{\max}=2\arcsin e},
\qquad
\boxed{P_{\max}=(0,\pm b)}.
\]

\textbf{最快理解：}角平分线的垂直方向上，焦点连线的投影为
$2a\sin(\theta/2)$，不能超过焦距 $2c$，所以 $\sin(\theta/2)\le c/a$。

\textbf{最快代数证明：}
\[
\sin^2\frac\theta2
=1-\frac{b^2}{PF_1\cdot PF_2}
\le1-\frac{b^2}{a^2}=e^2,
\]
原因是 $PF_1\cdot PF_2\le a^2$，等号恰在 $PF_1=PF_2=a$ 时成立。

\textbf{高考三种直接应用：}
\begin{enumerate}
    \item \textbf{已知 $e$ 求最大角：}直接用 $\theta_{\max}=2\arcsin e$。
    \item \textbf{存在 $\theta\ge\theta_0$：}等价于
    $e\ge\sin(\theta_0/2)$（$0\le\theta_0<\pi$）。
    \item \textbf{存在直角／钝角：}分别对应
    $e\ge\sqrt2/2$ 与 $e>\sqrt2/2$。
\end{enumerate}

\textbf{考场提醒：}不要把存在性误当成任意性；不要丢失
$0<e<1$；不能将本结论直接搬到双曲线。
\end{summarybox}
